\documentclass[a4paper,11pt]{article}
\usepackage{lecturenotes}
\usepackage{school-notes-extension}
\geometry{margin=18mm,top=18mm,bottom=18mm,includehead,
  headheight=24pt,headsep=6mm}
\hypersetup{pdftitle={Typesetting Scripts - Block Reference},
  pdfauthor={Typesetting Scripts},pdfsubject={lecturenotes v0.2.0 usage reference}}
\fancyhf{}
\fancyhead[L]{\sffamily\small\bfseries\color{NoteAccent}TYPESETTING SCRIPTS}
\fancyhead[R]{\sffamily\small\color{NoteMuted}Block reference}
\fancyfoot[L]{\sffamily\scriptsize\color{NoteMuted}lecturenotes v0.2.0}
\fancyfoot[R]{\sffamily\scriptsize\color{NoteMuted}\thepage}
\lstdefinestyle{referencecall}{language=[LaTeX]TeX,
  basicstyle=\fontsize{9}{11}\selectfont\ttfamily\color{NoteInk},
  keywordstyle=\color{NoteAccent},commentstyle=\color{NoteMuted},
  backgroundcolor=\color{NoteCode},frame=single,
  rulecolor=\color{NoteRule},framerule=0.3pt,framesep=7pt,
  columns=fullflexible,keepspaces=true,showstringspaces=false,
  breaklines=true,breakatwhitespace=false,numbers=none,
  aboveskip=5pt,belowskip=8pt,tabsize=2}
\lstset{style=referencecall}
\ExplSyntaxOn
\NewDocumentCommand{\code}{m}{
  \group_begin:
  \ttfamily
  \tl_set:Nx \l_tmpa_tl {\tl_to_str:n {#1}}
  % Remove the space detokenize adds before a command's argument delimiter.
  \regex_replace_all:nnN {(\\[A-Za-z@]+)\s+([\{\[])} {\1\2} \l_tmpa_tl
  \tl_trim_spaces:N \l_tmpa_tl
  \tl_use:N \l_tmpa_tl
  \group_end:
}
\ExplSyntaxOff
\newcommand{\tagline}[1]{\par\smallskip
  {\sffamily\scriptsize\bfseries\color{NoteAccent}#1}\par\nobreak}
\newcommand{\referencepage}[3]{\clearpage
  \pdfbookmark[0]{#2}{manual:#1}\phantomsection\label{ref:#1}
  {\sffamily\fontsize{21}{25}\selectfont\bfseries #2\par}
  {\sffamily\small\color{NoteMuted}#3\par}\medskip}
\newenvironment{result}{\tagline{RENDERED RESULT}
  \begin{tcolorbox}[enhanced,breakable,colback=white,colframe=NoteRule,
    boxrule=0.4pt,arc=0pt,left=9pt,right=9pt,top=6pt,bottom=6pt,
    before skip=4pt,after skip=8pt]}{\end{tcolorbox}}
\newcommand{\callfile}[1]{\tagline{CALL}
  \lstinputlisting{docs/reference/examples/#1.tex}}
\newcommand{\resultfile}[1]{\begin{result}
  \input{docs/reference/examples/#1.tex}\end{result}}
\newcommand{\example}[1]{\callfile{#1}\resultfile{#1}}
\newcommand{\indexrow}[3]{\hyperref[ref:#1]{#2} & #3 & \pageref{ref:#1}\\}
\NewDocumentEnvironment{options}{+b}{\tagline{OPTIONS AND DEFAULTS}
  \begingroup\small\renewcommand{\arraystretch}{1.2}
  \begin{tabularx}{\linewidth}{@{}p{54mm}X@{}}
    #1
  \end{tabularx}\endgroup\par\smallskip}{}
\begin{document}
\pdfbookmark[0]{Find a block}{manual:index}
{\sffamily\fontsize{27}{31}\selectfont\bfseries Block reference\par}
{\sffamily\large\color{NoteAccent}Calls, options, and actual output\par}
\medskip
Use this while writing a lecture. The gray area is code to copy; the bordered
area shows the result produced by the current package. Click a topic below or
open this PDF's bookmarks sidebar to jump to its reference page.

\tagline{FIND THE RIGHT BLOCK}
\begin{tabularx}{\linewidth}{@{}p{39mm}Xr@{}}
\textbf{Task} & \textbf{Command or environment} & \textbf{Page}\\\midrule
\indexrow{headings}{Organize topics}{\code{\noteheading}; heading aliases}
\indexrow{navigation}{Navigate the PDF}{\code{\notesetup}; contents and bookmarks}
\indexrow{text}{Explain or define}{\code{\notebox}, \code{\notedefinition}, \code{\assumption}}
\indexrow{equations}{Highlight a result}{\code{\keyequation}}
\indexrow{steps}{Explain a derivation}{\code{notederivation}, \code{\derivestep}}
\indexrow{legacy}{Align compact steps}{Legacy and multiline derivations}
\indexrow{tables}{List values or symbols}{\code{notetable}, \code{notation}}
\indexrow{panels}{Place content side by side}{\code{notepanels}, \code{notepanel}}
\indexrow{figures}{Insert an existing image}{\code{\notefigure}}
\indexrow{circuits}{Draw a circuit}{\code{notecircuit}; presets and annotations}
\indexrow{signals}{Plot a signal}{\code{notesignal}}
\indexrow{code}{Include source code}{\code{notecode}}
\indexrow{summary}{Summarize a topic}{\code{\notesummary} (School-Notes only)}
\indexrow{cli}{Check or update lectures}{\code{notes check}, \code{notes update-style}}
\indexrow{shortcuts}{Type blocks quickly}{VS Code shortcuts and common mistakes}
\bottomrule
\end{tabularx}

\tagline{FIRST DOCUMENT}
Place \code{lecturenotes.sty} beside \code{lecture.tex}. All block examples go
between \code{\begin{document}} and \code{\end{document}}.
\begin{lstlisting}
\documentclass[a4paper,11pt]{article}
\usepackage{lecturenotes}
\lectureheader{ELE734}{1}
\begin{document}
\tableofcontents
\clearpage
\noteheading{1}{First topic}
Your notes go here.
\end{document}
\end{lstlisting}
\textbf{Reading the syntax.} Braces \code{{...}} hold required arguments;
brackets \code{[...]} hold optional arguments. An environment needs matching
\code{\begin} and \code{\end}. In option lists, separate settings with commas
and brace values that contain commas. Compile twice after changing headings
or references. Existing lectures need the v0.2.0 style for the new blocks.

\referencepage{headings}{Headings}{\code{\noteheading[options]{level}{title}}}
Use levels 1 through 5 to express increasingly specific topics. The heading
itself is unnumbered. Start with level 1 and move down one level at a time.
\example{headings}
\begin{options}
\code{level} & Required integer from 1 to 5. Other levels cause an error.\\
\code{bookmark={...}} & Optional plain-text sidebar title. The printed heading
keeps its mathematics; the sidebar gets a readable title.\\
\end{options}
\tagline{EQUIVALENT OLDER NAMES}
\begin{lstlisting}
\notesection{Title}        % same as \noteheading{1}{Title}
\subnotesection{Title}     % same as \noteheading{2}{Title}
\subsubnotesection{Title}  % same as \noteheading{3}{Title}
\end{lstlisting}
Use \code{\noteheading{4}} and \code{\noteheading{5}} for deeper headings;
there is no need to repeat ``sub'' in the command name.

\textbf{Editor shortcuts:} \code{;h1} through \code{;h5}; \code{;sec} and
\code{;sub} for the older names. \code{;hbm} asks for the bookmark title,
then the level, then the printed title.

\referencepage{navigation}{Contents and PDF navigation}{\code{\notesetup{toc-depth=...,bookmark-depth=...}}}
\textbf{Contents} is a page printed inside your lecture PDF.
\textbf{Bookmarks} are a topic tree in the PDF viewer's sidebar.
Clicking either kind of entry jumps to the heading. These settings control
the two lists independently; every heading still appears in the lecture body.
\tagline{CALL: TWO LEVELS ON PAPER, FIVE IN THE SIDEBAR}
\begin{lstlisting}
% Set this before \tableofcontents and before your headings.
\notesetup{toc-depth=2,bookmark-depth=5}
\tableofcontents
\clearpage
% Now add the five headings shown on the previous page.
\end{lstlisting}
\tagline{PRINTED CONTENTS: ONLY LEVELS 1 AND 2}
\begin{result}
  \notesetup{toc-depth=2}
  \tableofcontents
\end{result}
\notesetup{toc-depth=5}
\tagline{SIDEBAR TREE FOR THE SAME FIVE HEADINGS}
\begin{lstlisting}[language={},basicstyle=\small\ttfamily]
Current mirrors
  Cascode mirrors
    Output resistance ro
      Dominant term
        Approximation condition
\end{lstlisting}
\begin{options}
\code{toc-depth=5} & Default. Values 0--5; 0 hides all note-heading entries
from the printed contents. You still need \code{\tableofcontents} to print it.\\
\code{bookmark-depth=5} & Default. Values 0--5; 0 hides note-heading entries
from the viewer's sidebar. It does not remove body headings.\\
\code{circuit-preset=none} & Another \code{\notesetup} setting. See page
\pageref{ref:circuits} for \code{compact} and \code{standard}.\\
\end{options}
\textbf{When the list looks stale:} compile twice. Open your PDF viewer's
bookmarks or outline sidebar to see bookmarks. The optional
heading \code{bookmark} text changes the sidebar title; the printed contents
keeps the printed heading text. \textbf{Shortcut:} \code{;nav}.

\referencepage{text}{Text, definitions, and assumptions}{Three blocks for three kinds of prose}
\textbf{General explanation:} \code{\notebox{title}{body}}.
Use for a physical interpretation, a reminder, or a grouped explanation.
\example{box}
\textbf{A term and its meaning:} \code{\notedefinition{term}{body}}.
The definition label and slim rule distinguish it from a general note.
\example{definition}
\textbf{A condition on the model:} \code{\assumption{body}}.
This stays inline rather than taking the space of another box.
\example{assumption}
Bodies accept normal LaTeX prose, inline mathematics such as \code{$\tau=RC$},
and paragraphs. The boxes can continue onto another page. These three commands
have no block-specific optional arguments.

\textbf{Shortcut:} \code{;def} inserts a definition. There are currently no
dedicated block shortcuts for \code{\notebox} or \code{\assumption}.

\referencepage{equations}{Key equations}{\code{\keyequation{title}{math}[explanation]}}
Use for a result worth finding again. The equation is displayed in a tinted
box, with an optional prose explanation below it.
\example{equation}
\textbf{Argument order matters:} the optional explanation comes
\emph{after} the title and equation. Omitting it gives the original
two-argument call. The equation argument is already in math mode: do not
surround it with \code{$...$}, \code{\[...\]}, or an \code{equation} environment.

\tagline{MORE THAN ONE EQUATION LINE}
Use \code{aligned} inside the equation argument. Put \code{&} where the
lines should align, and \code{\\} at each line break.
\example{equation-multiline}
\textbf{Numbering:} key equations are unnumbered. For an equation that needs
a numbered cross-reference, use a normal LaTeX \code{equation} environment
and \code{\label{eq:unique-name}}.

\textbf{Shortcut:} \code{;keq} inserts the title and equation. Add the final
\code{[explanation]} yourself when needed.

\referencepage{steps}{Derivations with explanations}{\code{notederivation} containing \code{\derivestep{math}{prose}}}
Use a step for each equation and the explanation that belongs to it.
Explanations wrap as ordinary prose. Each step stays together; the full
derivation can break between steps.
\example{derivation}
\tagline{PUT THE EXPLANATION BESIDE THE EQUATION}
\example{derivation-beside}
\begin{options}
\code{mode=auto} & Default. Detects a literal \code{\derivestep} in the body.
Use \code{mode=steps} when using \code{\input} or to state the layout explicitly.\\
\code{explanation=below} & Default. Choose \code{beside} for a text column on
the right. \code{explanations} is an accepted alias.\\
\code{equation-width} & Default \code{0.52\linewidth}; left-column width in
\code{beside} mode. Leave room for the explanation and the gap.\\
\code{gap=6mm} & Default. Space between columns in \code{beside} mode.\\
\end{options}
\textbf{Shortcuts:} \code{;drv} for the whole block; \code{;step} for another
step. The first math argument is already in math mode; the second is prose.

\referencepage{legacy}{Compact and multiline derivations}{The original alignment remains supported}
\textbf{Compact explanations.} Use the original \code{alignedat} body for
short annotations that fit on one line. The initial \code{&} starts the
equation column; \code{&&} starts its annotation column.
\example{derivation-legacy}
\code{mode=legacy} states the layout explicitly. Omitting the option also
selects it when the body has no literal \code{\derivestep}.
Annotations use \code{\text{...}} and do not wrap like prose. For longer
explanations, use the step layout on page \pageref{ref:steps}.

\tagline{MULTIPLE MATH LINES IN ONE WRAPPING STEP}
Inside \code{\derivestep}, \code{&} aligns equal signs and \code{\\}
starts another math line. Do not add another \code{aligned} wrapper;
the step already supplies one.
\example{derivation-multiline}
\textbf{Loading steps from a separate file:}
\begin{lstlisting}
\begin{notederivation}[mode=steps]{Voltage gain}
  \input{gain-steps.tex} % the file contains \derivestep calls
\end{notederivation}
\end{lstlisting}
\textbf{Shortcut:} \code{;der} inserts the original aligned-column layout.
Use \code{;drv} for wrapping explanations.

\referencepage{tables}{Tables and notation}{\code{notetable} for custom columns; \code{notation} for symbols}
\code{\begin{notetable}{column specification}{header cells}} starts a
full-width table. Use \code{l}, \code{c}, or \code{r} for natural-width columns;
\code{X} makes a paragraph column that fills the remaining width.
\example{table}
Separate cells with \code{&}. End each body row with \code{\\} or
\code{\tabularnewline}. The header's row ending and the horizontal rules
are automatic: do not add a trailing \code{\\} to the header argument.
The number of cells must match the column specification.

\tagline{A SYMBOL LIST WITH FIXED HEADERS}
\code{notation} supplies three columns: Symbol, Meaning, and Unit.
The middle column wraps. Write three cells per row.
\example{notation}
Tables do not break across pages. Keep long lists in separate tables.
For a literal ampersand in prose, write \code{\&} so it is not treated as
a new cell.

\textbf{Shortcut:} \code{;tbl} inserts a custom notes table. There is no
dedicated \code{notation} shortcut.

\referencepage{panels}{Content side by side}{\code{notepanels} contains two or three \code{notepanel} columns}
Put a diagram beside its explanation, compare two models, or show three
stages. Panel titles are required. Widths are relative weights, not lengths.
\example{panels}
\tagline{THREE EQUAL COLUMNS}
\example{panels-three}
\begin{options}
\code{columns=2} & Default. Only 2 or 3 is supported. Additional panels
start another row after this many columns.\\
\code{widths={}} & Default: equal widths. \code{widths={2,1}} gives the first
panel twice the usable width. Supply one positive number per column.\\
\code{gap=6mm} & Default. Gaps are removed before widths are distributed.\\
\code{align=top} & Default. Also \code{center} or \code{bottom}.\\
\code{keep-together=true} & Default. Keeps the entire block on one page.
\code{false} permits breaks between complete rows.\\
\end{options}
Each panel is a minipage and cannot split internally. Use inline figures
(\code{placement=here}) inside it; normal floats cannot go inside a minipage.
A panel must be inside \code{notepanels}. Leave enough space for the whole row.
\textbf{Shortcuts:} \code{;pan2}, \code{;pan3}, and \code{;panel}.

\referencepage{figures}{Images with captions}{\code{\notefigure[options]{image file}{caption}}}
Use for an existing PNG, JPEG, or PDF figure. It fits within both dimension
limits, preserves its aspect ratio, and adds a numbered caption.
\example{figure}
\begin{options}
\code{width=0.85\linewidth} & Default maximum width. \code{\linewidth}
adapts to the current panel or text area.\\
\code{maxheight=65mm} & Default maximum height. The smaller of the two bounds
determines the final scale; the image is not stretched.\\
\code{placement=here} & Default. Image and caption stay together, inline at
the call. Use this inside panels.\\
\code{placement=float} & Uses normal \code{htbp} floating. Placement letters
such as \code{t} or \code{htbp!} are also supported; LaTeX chooses the final position.\\
\code{label={}} & Default: no reference target. Set a unique name such as
\code{fig:rc-response} and use \code{\ref{fig:rc-response}}.\\
\end{options}
The example's long path is for this reference's included asset. In a lecture,
put the image in \code{figures/} and use \code{{figures/rc-response.pdf}}.
Paths are relative to the document's build directory. You may set
\code{\graphicspath{{figures/}}} and use just the file name.
Both dimensions must be positive. Compile twice to resolve figure numbers.

\textbf{Shortcut:} \code{;fig} asks for the file name, caption, and label
suffix. It supplies \code{figures/}, \code{fig:}, width \code{0.8\linewidth},
and inline placement. Replace every field with your actual values.

\referencepage{circuits}{Circuits and annotations}{\code{\begin{notecircuit}[drawing options]{title}}}
Write ordinary CircuitikZ commands inside the block. Set a preset before the
diagram to avoid repeating the common scale and symbol options.
\example{circuit}
\begin{options}
\code{circuit-preset=none} & Default. Supplies no drawing options; existing
per-diagram settings remain in control.\\
\code{circuit-preset=compact} & American symbols, scale 0.55, transformed
labels, and compact current arrows. A useful starting point for notes.\\
\code{circuit-preset=standard} & Same conventions at scale 0.8.\\
\code{note-highlight} & TikZ style: translucent yellow, rounded, 5pt stroke.
Draw it before the dark wire so the connection remains clear.\\
\code{note-current} & TikZ style: thin red current arrow with a small tip.
Put its label in a separate node with a little space from the arrow.\\
\end{options}
\textbf{Orthogonal MOS connection.} \code{(M.G |- bias)} uses the gate's
horizontal coordinate and the bias node's vertical coordinate. Connecting
through that point gives horizontal and vertical wires even when the symbol
is mirrored. It avoids a guessed gate coordinate and an angled wire.

Explicit drawing options override the preset, for example
\code{\begin{notecircuit}[scale=0.65]{Title}}. The block centers the diagram
and adds its title; it does not number a figure or create a floating caption.
Keep wire runs short and give labels space.
\textbf{Shortcuts:} \code{;cset}, \code{;cir}, \code{;cur}, \code{;mark}.

\referencepage{signals}{Signal and response plots}{\code{\begin{notesignal}[axis options]{title}}}
Use for a mathematical curve or a set of data points. The block provides
a centered PGFPlots axis, muted grid, and teal plot strokes.
\example{signal}
\tagline{OPTIONS YOU WILL USUALLY CHANGE}
\begin{options}
\code{xlabel}, \code{ylabel} & Axis labels. Brace values containing math,
spaces, or commas, as shown above.\\
\code{xmin}, \code{xmax} & Visible horizontal range. Use \code{ymin} and
\code{ymax} for the vertical range.\\
\code{xtick}, \code{ytick} & Explicit tick lists, for example \code{xtick={0,1,2}}.\\
\code{width}, \code{height} & Defaults \code{0.94\linewidth} and \code{5cm}.
The example uses a slightly shorter axis.\\
\code{domain}, \code{samples} & Options on \code{\addplot} for a formula,
not required arguments of the block.\\
\end{options}
The optional list accepts normal PGFPlots axis settings. Use PGFPlots
expression syntax inside a formula: \code{1-exp(-x)}, not a LaTeX equation.
For data, use \code{\addplot coordinates {(0,0) (1,0.5) (2,0.8)};}.
Add another \code{\addplot} for another curve and \code{\addlegendentry{...}}
after a plot for a legend entry. End drawing commands with semicolons.
There is no dedicated \code{notesignal} shortcut.

\referencepage{code}{Source-code blocks}{\code{\begin{notecode}[listings options]}}
Code is read verbatim: type underscores, braces, and \texttt{\#} normally.
The block supplies a neutral background, thin frame, syntax colors, and
automatic line wrapping.
\example{code-c}
\tagline{ARM ASSEMBLY}
\example{code-arm}
\begin{options}
\code{language=C} & Default. Choose another language supported by
\code{listings}, or the bundled \code{ARM} lexer.\\
\code{numbers=none} & Default. \code{numbers=left} adds line numbers.\\
\code{firstnumber=1} & Standard listings option; change the starting line
number when presenting a later part of a program.\\
\end{options}
The ARM lexer highlights a small core of instructions and \code{;} comments;
it is a display helper, not an assembler or validator.
Place \code{notecode} directly in the document. Verbatim environments do not
work inside a macro argument such as the body of \code{\notebox}, or inside
the captured body of \code{notepanels}. Keep the closing
\code{\end{notecode}} on its own line. There is no dedicated block shortcut.

\referencepage{summary}{School-Notes summary block}{\code{\notesummary{title}{overview}{table rows}} - local extension}
This command exists in the updated lecture styles in your School-Notes
directory. It is preserved when updating those styles, but is \textbf{not}
bundled into the style generated by \code{notes new}.
\example{summary}
\tagline{WHAT EACH ARGUMENT DOES}
\begin{options}
\code{title} & Required title at the top of the white note box.\\
\code{overview} & Required prose in that box. Inline math and paragraphs
work like normal LaTeX.\\
\code{table rows} & Required two-column rows below the box. Use a short
label, \code{&}, the description, then \code{\\}.\\
\end{options}
The label column is fixed at \code{4em}; the description column fills the
remaining width and wraps. Keep labels short, such as Model, Limit, or Use.
There are no automatically inserted table headings and no optional arguments.

\tagline{IF THE COMMAND IS UNDEFINED IN A NEW LECTURE}
Use \code{\notebox} followed by \code{notetable} for the same general
structure, or intentionally copy the local definition into that lecture's
style. The reference loads a separate copy of this extension only to show
its output; the canonical package remains the source for every other example.

\textbf{Editor note:} \code{;sum} inserts a mathematical summation, not
\code{\notesummary}. There is no summary-block shortcut at present.

\referencepage{cli}{Checking and updating lectures}{Terminal commands that support the writing workflow}
\tagline{CHECK A LECTURE, COURSE, OR THE WHOLE NOTES ROOT}
\begin{lstlisting}[language={}]
notes check ELE734/lecture-01
notes check ELE734 --compile
notes check . --strict
\end{lstlisting}
Run commands from your School-Notes root, or supply a path to it. Omitting
the path checks the current directory. The static checker follows literal
\code{\input} and \code{\include} files and reports duplicate labels,
undefined references, missing files, placeholder labels, and unfinished notes.
\tagline{ILLUSTRATIVE OUTPUT}
\begin{lstlisting}[language={}]
ELE734/lecture-01/lecture.tex:42: error [undefined-reference]:
  reference "fig:gain" has no label
\end{lstlisting}
\begin{options}
\code{--compile} & Runs pdfLaTeX twice without shell escape, in a temporary
output directory. Reports build errors and layout warnings while preserving
existing PDFs, logs, and auxiliary files. Each pass has a 45-second timeout.\\
\code{--strict} & Warnings also produce exit status 1. Normally warnings
return 0, errors return 1, and invalid command syntax returns 2.\\
\end{options}
The static scan does not expand arbitrary TeX macros or know every
package-generated label. Use the real build to assess those cases.

\tagline{PREVIEW, THEN APPLY A STYLE UPDATE}
\begin{lstlisting}[language={}]
notes update-style ELE734 --dry-run
notes update-style ELE734
\end{lstlisting}
This accepts a lecture, course, or notes root; omission means the current
directory. The preview reports versions, custom edits, and conflicts. Applying
merges compatible local edits into the current bundled style, including local
colors and added commands such as \code{\notesummary}.

\textbf{Concrete file change:} \code{lecturenotes.sty} becomes the updated
style. Its exact previous bytes go into \code{lecturenotes.sty.bak}, then
\code{.bak.1}, \code{.bak.2}, and so on if earlier backups exist. Your
\code{lecture.tex} and figures are unchanged. Already-current styles get no
new backup. Rebuild the lecture afterward to see the updated formatting.

\textbf{Conflicts:} all lectures are checked before any style is replaced.
Inspect the reported conflict first. \code{--resolve=keep} favors conflicting
local edits; \code{--resolve=replace} favors the updated definition. Both
retain compatible additions. Keeping an old definition can retain old behavior.
Unknown versions and complex or structurally invalid merges need a manual merge.

\referencepage{shortcuts}{Editor shortcuts and quick fixes}{Keep this page open beside VS Code}
In your School-Notes workspace, type these at the start of a line outside
math mode. HyperSnips expands the block; use Tab to move through its fields.
After updating snippets, run \textbf{HyperSnips: Reload Snippets} from the
Command Palette. These shortcuts belong to the School-Notes workspace.
\begin{notetable}{lX}{\textbf{Shortcut} & \textbf{Inserts}}
\code{;h1}--\code{;h5} & Heading at the chosen level.\\
\code{;sec}, \code{;sub} & Main and second-level heading aliases.\\
\code{;hbm} & Heading with plain bookmark title; fields: bookmark, level, title.\\
\code{;nav} & Printed-contents depth and sidebar depth.\\
\code{;def} & Definition with a term and body.\\
\code{;keq} & Key equation with title and math.\\
\code{;der} & Original derivation with short aligned annotations.\\
\code{;drv}, \code{;step} & Wrapping derivation and an additional step.\\
\code{;tbl} & Table specification, header, and body row.\\
\code{;pan2}, \code{;pan3} & Two or three titled panels.\\
\code{;panel} & Additional panel inside \code{notepanels}.\\
\code{;fig} & Image, caption, and label; fields: file name, caption, label suffix.\\
\code{;cset}, \code{;cir} & Circuit preset and circuit block.\\
\code{;cur}, \code{;mark} & Current arrow and highlighted connection.\\
\end{notetable}
\tagline{COMMON MISTAKES}
\begin{notetable}{>{\raggedright\arraybackslash}p{43mm}X}{\textbf{Symptom} & \textbf{Fix}}
Undefined new command & Update that lecture's local style to v0.2.0.
For \code{\notesummary}, see page \pageref{ref:summary}.\\
An equation gives a math error & Remove the outer \code{$...$} or
\code{\[...\]} from \code{\keyequation} or \code{\derivestep}.\\
Extra alignment tab error & Match each table row's number of cells to
the columns. Escape literal ampersands as \code{\&}.\\
Figure number is ``??'' & Give the figure a unique label, use exactly the
same name in \code{\ref}, and compile twice.\\
A panel or step runs off the page & Shorten or split its content. A panel
or individual step cannot split internally.\\
Floating figure fails in a panel & Set \code{placement=here} for that figure.\\
Derivation loaded by \code{\input} fails & Set \code{mode=steps} on
\code{notederivation}.\\
\end{notetable}
\textbf{Copyable source:} the illustrated block examples live in
\code{docs/reference/examples/}. The reference includes those same files to
render the results, so the calls and outputs stay in sync.
\end{document}
